
\chapter{光的干涉}

\begin{definition}[可见光]
可见光的波长范围：400 nm到760 nm，其中人眼最为敏感的是波长为555nm的黄绿光。
\end{definition}

光的本性经历过一条很长的认识路线。

\textbf{波动性证据链}大致是这样的：
\begin{itemize}
 \item 牛顿曾提出微粒说，认为光是由光源发出的微粒流；
 \item 惠更斯提出波动说，认为光是一种波；
 \item 杨氏双缝实验直接观察到稳定干涉条纹，强有力地支持了光的波动性；
 \item 麦克斯韦进一步指出，光本质上是一定频率范围内的电磁波。
\end{itemize}
光是电磁波，其中电场 \(\vb*{E}\) 和磁场 \(\vb*{H}\) 垂直于传播方向且互相垂直。
\begin{itemize}
 \item 光波中有一个振荡的电场 \(\vb*{E}\)，它的振动方向就是光的"偏振方向"
 \item  光的强度（即能量流密度）正比于电场振幅的平方：\(I\propto E_m^2\)；
 \item 探测器接收到的是时间平均值 \(I\propto\langle E^2\rangle\)，因为光的振荡频率远高于任何探测器的响应速度。
\end{itemize}

\begin{definition}[光矢量，光振动，光波的振动方向]
若光沿 \(x\) 方向传播，一个最简形式的光波可以写成
\[
E=E_m\cos\omega\left(t-\frac{x}{u}\right),\qquad
H=H_m\cos\omega\left(t-\frac{x}{u}\right).
\]
在波动光学中，通常把光波中的电场强度矢量 \(\vb*{E}\) 称为\textbf{光矢量}。\\
光波中电场强度矢量随时间作周期性变化，称为\textbf{光振动}。\\
光矢量 \(\vb*{E}\) 的振动方向，称为\textbf{光波的振动方向}。
\end{definition}

\begin{note}
 题目只要说“光的振动方向”“偏振方向”，默认指的都是电场 \(\vb*{E}\) 的方向； 它把后面偏振题里真正会被筛选、会被比较夹角的那个方向固定下来； 传播方向、\(\vb*{E}\)、\(\vb*{H}\) 两两垂直，而偏振片筛的是 \(\vb*{E}\) 在某个横向方向上的分量； 把“光沿哪边传播”和“电矢量朝哪边振动”混成同一个方向。
\end{note}

\begin{definition}[相干光]\label{def:coherent-light}
频率相同、振动方向相同、相位差恒定的两列光，称为\textbf{相干光}。
\begin{itemize}
 \item \textbf{频率相同：} 保证两束光不会越走越“拍”，否则条纹位置会不停漂；
 \item \textbf{振动方向相同：} 保证它们真能在同一方向上发生叠加；
 \item \textbf{相位差恒定：} 保证明暗分布稳定，不会一会儿亮一会儿暗。
\end{itemize}
注意：没有要求光的强度。
\end{definition}

\begin{definition}[普通光源]
除激光源之外，由大量原子、分子或离子发光形成的常见光源，都可以统称为\textbf{普通光源}。普通光源最典型的发光机制是\textbf{自发辐射}：原子先被激发到高能级，随后自发跃迁回较低能级，同时放出一段光波。普通光源发出的并不是“无限长的连续正弦波”，而更像是一段一段有限长度的光波列。若平均发光持续时间为 \(\Delta t\)，则这段光波列的长度大约为 \(L=c\Delta t\)。普通光源中这个时间通常非常短，因此波列长度也有限。普通光源天然不相干:
\begin{itemize}
 \item \textbf{间歇性：} 各原子并不是整齐划一地一直发光，而是断断续续地发出有限长度的波列；所以不同原子发出的两束光通常是独立的；
 \item \textbf{随机性：} 每次发光的振动方向和初相都带有随机性，不同原子之间更没有固定相位关系。所以同一原子先后两次发出的光也通常是独立的。
\end{itemize}
“普通光源不相干”并不是说它不能参与干涉，而是说\textbf{不能直接把两束彼此独立的普通光拿来就稳定干涉。} 若想看到稳定条纹，先抓住同一原子某一次自发辐射形成的那一小段波列，再把这\textbf{同一波列}拆成两路，让它们走不同光路后重新相遇。 \[\text{同一波列分成两路} \implies \text{更容易获得稳定相位差} \implies \text{可形成相干光}.\]获得相干光最常见的两条路径是：
\begin{itemize}
 \item \textbf{分波阵面法：} 从同一波前上截取两部分光，让它们传播后重新叠加。它更强调同一波前不同区域之间的\textbf{空间相干}。
 \item \textbf{分振幅法：} 让同一束光在界面上部分反射、部分透射，分成两束后再相遇。它更强调同一波列前后部分之间的\textbf{时间相干}。
\end{itemize}
\end{definition}

\begin{center}
\includegraphics[width=0.6\linewidth]{pdf-exercises/3.png}
\end{center}
